% ==============================================================================
% MODULE 08: CONCLUSION AND FUTURE SCOPE
% ==============================================================================

\chapter{Conclusion and Future Scope}
\label{chap:conclusion}

Traditional video streaming platforms store massive video libraries in centralized cloud data centers, creating high recurring costs for storage and bandwidth. \textbf{CustardFin} solves this problem by keeping video files on the owner's home computer and using a lightweight cloud server only for catalog search, connection coordination, and temporary caching. This approach provides reliable video streaming, protects personal privacy, and reduces cloud operating costs by 95.89\%.

% ------------------------------------------------------------------------------
\section{Summary of Achievements}
\label{sec:conclusion_achievements}
% ------------------------------------------------------------------------------

CustardFin successfully satisfies all core functional requirements (FR-1 through FR-10) and fulfills all ten non-negotiable architectural invariants:

\begin{itemize}
    \item \textbf{Zero Cloud Video Storage (Invariant~1):} Original video files remain permanently on home drives. The cloud stores only 650-byte JSON metadata manifests, eliminating perpetual cloud storage bills.
    \item \textbf{Home CGNAT Traversal (Invariant~2):} Home nodes connect to the cloud using outbound-only TLS reverse tunnels (FRP). Users never need to open router ports, configure port forwarding, or pay for static public IP addresses.
    \item \textbf{100\% Free and Open-Source Stack (Invariant~3):} The system is built entirely on open-source technologies (Go, PostgreSQL, Redis, MinIO/Floci, FRP, Next.js, and SQLite) without proprietary software or licensing fees.
    \item \textbf{High-Speed Hardware Transcoding (Invariant~6):} Pre-encoded HLS with NVIDIA NVENC delivers $10.74\times$ real-time encoding speed ($322+\text{ fps}$) while keeping CPU usage under $5\%$, supported by automatic fallback to CPU software encoding.
    \item \textbf{Origin Shield Bandwidth Protection (Invariant~7):} The 24-hour pull-through cache uploads hot segments once and serves all subsequent viewers from the cache. In a 10-viewer test, the cache offloaded $90\%$ of traffic, allowing a modest $20\text{ Mbps}$ home connection to serve 10 viewers smoothly.
    \item \textbf{Data Privacy and Sovereignty (Invariant~8):} Private watch history, resume bookmarks, and unshared home videos remain strictly on the local edge computer in its SQLite database.
    \item \textbf{Secure Signed URLs (Invariant~9):} Video chunks are protected by short-lived HMAC-SHA256 signed tokens, preventing hotlinking and unauthorized access.
    \item \textbf{Admission Control and Capacity Limits (Invariant~10):} A limit of 10 concurrent remote streams per node protects home bandwidth from overload, returning clear \texttt{HTTP 429 Too Many Requests} responses when full.
    \item \textbf{Fast LAN and Offline Streaming:} Local streaming operates with sub-4ms latency, and the embedded Mini-Jellyfin player provides full playback even during internet or cloud outages.
    \item \textbf{Verified 95.89\% Cloud Cost Reduction:} Monthly cloud operating costs drop from \rupee~1{,}00{,}380.00 (\$1{,}195.00) in traditional architectures to \rupee~4{,}128.60 (\$49.15) with CustardFin, saving over \rupee~34.65 Lakhs (\$41,250) over three years for a 100-user community.
\end{itemize}

% ------------------------------------------------------------------------------
\section{Current Limitations (Version 1)}
\label{sec:conclusion_limitations}
% ------------------------------------------------------------------------------

While CustardFin provides a fully functional streaming platform, several limitations exist in Version~1:
\begin{enumerate}
    \item \textbf{Single Relay Server Dependency:} Uncached chunks must travel through a single cloud relay server running \texttt{frps}. If this relay goes down, remote playback for uncached chunks is temporarily interrupted.
    \item \textbf{No Direct Browser P2P Swarming:} In Version~1, edge nodes stream through the cloud relay; direct peer-to-peer chunk sharing between browser viewers using WebRTC is not yet implemented.
    \item \textbf{No Proprietary DRM:} To remain completely free and open-source, CustardFin relies on HMAC-SHA256 signed URLs and HLS segmentation rather than commercial DRM systems like Google Widevine or Apple FairPlay.
    \item \textbf{Web-First Player:} The frontend is currently a responsive web application for desktop and mobile browsers. Dedicated native applications for smart TVs (Android TV, Apple TV) are not yet available.
\end{enumerate}

% ------------------------------------------------------------------------------
\section{Future Enhancements}
\label{sec:conclusion_future}
% ------------------------------------------------------------------------------

The modular design of CustardFin allows for several planned enhancements:
\begin{itemize}
    \item \textbf{Distributed Multi-Region Relays:} Deploying reverse tunnel relays across multiple cloud regions using GeoDNS will reduce latency for global viewers and eliminate single points of failure.
    \item \textbf{WebRTC Peer-Assisted Streaming:} Enabling viewers watching the same popular video to share cached chunks directly through browser WebRTC data channels will further reduce cloud bandwidth.
    \item \textbf{Cryptographic Manifest Verification:} Adding Ed25519 digital signatures to public manifests will verify node identity and metadata integrity.
    \item \textbf{Smart TV Applications:} Native apps built with Flutter or React Native can bring CustardFin to smart TVs with automatic local network discovery.
\end{itemize}

% ------------------------------------------------------------------------------
\section{Concluding Remarks}
\label{sec:conclusion_remarks}
% ------------------------------------------------------------------------------

CustardFin proves that high-quality media streaming does not require expensive, centralized cloud storage. By keeping video files safely on the creator's own computer and using a lightweight cloud server for discovery and temporary caching, CustardFin offers a private, community-focused, and highly economical streaming platform. This project successfully fulfills all academic and practical requirements of the MCA Cloud Computing Laboratory curriculum.
